\chapter{Fitting functions and emulators}\label{ch:fitting-emulators}
\chapterstatus{written}{2026-11-20}

Beside the halo models of \cref{ch:halo-model}, \fml{} has the standard fitting functions and emulators of the nonlinear
matter power spectrum: e-MANTIS, baccoemu (also its baryonic boosts and the Lagrangian bias expansion of tracers), Halofit
with the neutrino correction of Bird et al., EuclidEmulator2, the $f(R)$ and Brans--Dicke boosts, the Eisenstein--Hu transfer functions, and a general neural-network emulator that can be trained in C++ or loaded
from PyTorch. The solver's selectable nonlinear prescriptions are listed in \cref{tab:nonlinear-options}.

\begin{modulebox}{Cosmology/Spectra, Emulator}
\nsp{FML::COSMOLOGY}: \file{Cosmology/Spectra/Halofit.h} (\cls{Halofit}, versions \code{Takahashi2012} and \code{Bird}),
\file{EisensteinHu.h}, \file{PowerSpectrum.h} (the solver's nonlinear tables, \code{solve_jbd_boost}).\par
\nsp{FML::EMULATOR}: \file{Emulator/EuclidEmulator2.h} (\nsp{EUCLIDEMULATOR2}, data file \file{EuclidEmulator2.dat}),
\file{Emulator/EMantis/EMantis.h} (\nsp{EMANTIS}), \file{Emulator/Bacco/Bacco.h} (\nsp{BACCO}; networks fetched by
\file{get_data.sh}), \file{Emulator/FofrFittingFunction.h} (\nsp{FOFRFITTINGFUNCTION}), \file{Emulator/NeuralNetwork/NeuralNetwork.h}
(\nsp{NEURALNETWORK}) with \file{export_lightning_emulator.py}; the training points from \file{LatinHyperCube}.\par
Documentation: \file{Cosmology/Spectra/NONLINEAR.md}, \file{Cosmology/EinsteinBoltzmannSolver/OPTIONS.md},
\file{Emulator/NeuralNetwork/README.md}, \file{Emulator/EMantis/README.md}, \file{Emulator/Bacco/README.md}. Tests:
\file{tests/test_emulator.cpp}, \file{Emulator/EMantis/tests}, \file{Emulator/Bacco/tests}, the solver's \file{test_nonlinear},
\file{test_brans_dicke}, \file{Emulator/NeuralNetwork/tests}.
\end{modulebox}

\begin{table}[ht]
\centering\small
\caption{The \param{nonlinear} option of the Einstein--Boltzmann solver. The nonlinear total-matter power is the fourth column
of \file{matter_power*.txt}; with \code{lensing = true} the same ratio multiplies the Weyl power in the Limber part.}
\label{tab:nonlinear-options}
\begin{tabularx}{\textwidth}{@{}lX@{}}
\toprule
\param{nonlinear} & Method \\
\midrule
\code{none} & linear (default) \\
\code{halofit} & Takahashi et al.\ (2012) Halofit \cite{Takahashi2012} \\
\code{halofit_bird} & Halofit with the Bird et al.\ neutrino correction as in CLASS \cite{Bird2012} \\
\code{hmcode2016} & HMcode-2016 as CAMB \code{mead2016}, with neutrinos (\cref{sec:hm-hmcode2016}) \\
\code{hmcode2020} & gravity-only HMcode-2020, with neutrinos (\cref{sec:hm-hmcode2020}) \\
\code{react} & the halo-model reaction times HMcode-2020 of the pseudo cosmology (\cref{sec:hm-react}) \\
\code{emantis} & e-MANTIS $f(R)$ boost times a matched GR nonlinear spectrum (\cref{sec:fit-emantis}) \\
\code{ee2} & EuclidEmulator2 (flat \code{lcdm}, \code{w0wa}, within its ranges) \\
\code{jbd_boost} & Brans--Dicke: Sletmoen \& Winther (2024) from the mapped GR universe \\
\code{axionhmcode} & mixed cold/ultralight-axion halo model for \code{model = axion} (\cref{sec:hm-axion}) \\
\bottomrule
\end{tabularx}
\end{table}

\section{Halofit}\label{sec:fit-halofit}

Halofit \cite{Smith2003,Takahashi2012} writes $\Delta^2_{\rm NL}=\Delta^2_Q+\Delta^2_H$ in terms of $y=k/k_\sigma$, where
$k_\sigma$ is defined by a unit Gaussian-filtered variance, $\sigma^2(1/k_\sigma)=1$, and the effective slope and curvature
$n_{\rm eff}=-3-\dd\ln\sigma^2/\dd\ln R$, $C=-\dd^2\ln\sigma^2/\dd\ln R^2$ at that scale:
\begin{equation}
\Delta^2_Q=\Delta^2_L\frac{(1+\Delta^2_L)^{\beta_n}}{1+\alpha_n\Delta^2_L}e^{-y/4-y^2/8},\qquad
\Delta^2_H=\frac{a_ny^{3f_1}}{1+b_ny^{f_2}+(c_nf_3y)^{3-\gamma_n}}\,\frac{1}{1+\mu_ny^{-1}+\nu_ny^{-2}},
\end{equation}
with the coefficients $a_n,\ldots,\nu_n$ polynomial in $n_{\rm eff}$ and $C$ and with a $w$ dependence (Takahashi's
revision), and $f_{1,2,3}$ powers of $\Om(z)$. \code{Version::Bird} adds the massive-neutrino terms of Bird et al.\
\cite{Bird2012} in the form used by current CLASS (changes to $\beta_n$, the quasi-linear term and the one-halo amplitude;
the obsolete $\Om$-dependent amplitude is left out), with $f_\nu=\Omega_\nu/\Om$ today; for $f_\nu=0$ it is exactly
Takahashi. The class tabulates $\ln(P_{\rm NL}/P_L)$ in $(\ln a,\ln k)$ from $z=10$; the static \code{Halofit::nonlinear}
works on one redshift. Bird agrees with CLASS (classy 3.2.0) to 0.051\% at $z=0$--2, $k=0.01$--$10\hMpc$, with identical
linear input.

\section{EuclidEmulator2}\label{sec:fit-ee2}

EuclidEmulator2 \cite{Knabenhans2021} emulates the nonlinear boost $B(k,z)=P_{\rm NL}/P_L$ from a suite of
\textsc{PKDGRAV3} simulations in eight parameters $(\Omega_b,\Om,\sum m_\nu,n_s,h,w_0,w_a,A_s)$: the boost is decomposed by
principal components and each component expanded in a sparse polynomial chaos of Legendre polynomials (up to degree 16) of the
parameters, mapped to $[-1,1]$. The port in \file{EuclidEmulator2.h} keeps the original algorithm and data file, with SI
constants from \file{Units.h}, inline functions (several translation units), owned GSL objects, and errors (exceptions)
instead of crashes: outside $\Omega_b\in[0.04,0.06]$, $\Om\in[0.24,0.40]$, $\sum m_\nu\in[0,0.15]$\,eV, $n_s\in[0.92,1]$,
$h\in[0.61,0.73]$, $w_0\in[-1.3,-0.7]$, $w_a\in[-0.7,0.5]$, $A_s\in[1.7,2.5]\times10^{-9}$ it throws, as for more than 101
redshifts or a missing data file. The $k$ range is $0.0087$--$9.4\hMpc$ (the boost is held above and 1 below in the
solver); $z\le10$.

\begin{cppcode}[EuclidEmulator2 boost]
#include <FML/Emulator/EuclidEmulator2.h>
using namespace FML::EMULATOR::EUCLIDEMULATOR2;
EuclidEmulator ee2(0.049, 0.315, 0.06, 0.966, 0.674, -1.0, 0.0, 2.1e-9);
auto [k, B] = ee2.compute_boost(1.0);      // k [h/Mpc], B = P_NL / P_L at z = 1
\end{cppcode}

\section{e-MANTIS}\label{sec:fit-emantis}

\file{Emulator/EMantis/EMantis.h} ports the prediction part of
\href{https://gitlab.obspm.fr/e-mantis/e-mantis}{e-MANTIS 1.4.0}, including its trained Gaussian processes and
principal-component reconstruction. Cite the $f(R)$ power paper \cite{SaezCasares2023}, the halo-mass-function paper
\cite{SaezCasares2024}, and the software \cite{EMantisSoftware} as applicable. The wCDM power boost has no dedicated
paper in this release; cite the software. Training and cross-validation are not ported. e-MANTIS is distributed under
the GPL (version 3 or later), and so is the port, like \fml{} itself: \file{Emulator/EMantis} keeps the licence text
(\file{COPYING}) and the notices of the original.

\cls{NonLinearMGBoostEmulator} returns $P_{f(R)}/P_{\rm GR}$; \cls{NonLinearBoostEmulator} returns the wCDM
$P_{\rm NL}/P_L$; \cls{HMFEmulator} returns $\dd n/\dd\ln M$ for spherical-overdensity or FoF masses;
\cls{Sigma8Emulator} converts $A_s$ to $\sigma_8$. Units are $k$ in $h$/Mpc and mass in $M_\odot/h$.
Queries outside the trained ranges throw unless a documented extrapolation option is selected.
The data files are exported from the Python package; see \file{Emulator/EMantis/README.md} for the full ranges.

For \code{model = fr}, \code{fR_n = 1}, the solver's \code{nonlinear = emantis} computes
$P_{\rm NL,fR}=P_{\rm NL,GR} B_{fR}$ from a second GR perturbation solve with the same initial conditions.
\param{emantis_gr} chooses \code{ee2} (default), \code{hmcode2020}, \code{hmcode2016}, \code{halofit} or
\code{halofit_bird}. Flat geometry and a single perturbation component are required.
For nonzero $f_{R0}$ the allowed range is $10^{-7}\le |f_{R0}|\le10^{-4}$; zero gives GR.
Cosmological inputs outside the emulator range throw, or clamp with \param{emantis_extrapolate_cosmo}.
For $z\le3$ the boost is interpolated on \param{emantis_n_z} nodes; below the first bin
($0.029\hMpc$) it is extrapolated toward unity and above $9.7\hMpc$ it is held fixed.
At $z>3$ the solver continues the boost from $z=3$ with the linear $f(R)$/GR ratio.
SimulationCode also uses e-MANTIS for its $n=1$ $f(R)$ nonlinear comparison column, with an EE2 GR baseline;
there it clamps cosmological inputs and returns linear theory above $z=3$ (\file{SimulationCode/ANALYSIS.md}).
The older Winther fitting-function API remains available.

\begin{cppcode}[Nonlinear f(R) boost]
#include <FML/Emulator/EMantis/EMantis.h>
using namespace FML::EMULATOR::EMANTIS;
Parameters p;
p.Omega_b = 0.049; p.Omega_m = 0.31; p.h = 0.68; p.n_s = 0.965;
p.sigma8 = 0.81; p.logfR0 = 5.0;   // sigma8 of matched LCDM; |fR0| = 1e-5
NonLinearMGBoostEmulator emulator("fR");
auto B = emulator.boost(p, 1.0, {0.1, 1.0});
\end{cppcode}

\begin{validation}[Validation: e-MANTIS port]
\file{Emulator/EMantis/tests} compares with Python 1.4.0: typically $10^{-10}$ relative,
a few $10^{-8}$ when converting $A_s$ to $\sigma_8$. These measure port reproducibility,
not the percent-level physical emulator error. Solver boost-table interpolation is checked separately by
\file{test_emantis_boost}; extrapolated results have no additional calibration.
\end{validation}

\section{baccoemu}\label{sec:fit-bacco}

\file{Emulator/Bacco/Bacco.h} ports the prediction part of
\href{https://bitbucket.org/rangulo/baccoemu}{baccoemu 2.3.0} (MIT licence, kept in \file{Emulator/Bacco/LICENSE}):
neural-network emulators trained on the BACCO simulations, with massive neutrinos and $w_0$--$w_a$ dark energy. Cite the
papers of the parts used: the linear spectra \cite{Arico2021Linear}, the nonlinear boost \cite{Angulo2021} (model
\code{Angulo2021}, the default) or \cite{Arico2023} (\code{Arico2023}, wider ranges), the baryonic boosts
\cite{Arico2021Baryons} (\code{Arico2021}) and \cite{Burger2025} (\code{Burger2025}, also for the bispectrum), and the
Lagrangian bias expansion in real \cite{Zennaro2023} and redshift space \cite{PellejeroIbanez2023} (its LPT spectra are an
emulator of velocileptors \cite{Chen2021}).

The conventions are baccoemu's: $\Omega_{\rm cold}=\Omega_c+\Omega_b$, $\Omega_\nu=M_\nu/(93.14\,h^2)$, $\sigma_{8,\rm cold}$ of
the cold matter at $z=0$ (or $A_s$ at $0.05\,{\rm Mpc}^{-1}$, converted with the $\sigma_8$ emulator), $k$ in $\hMpc$.
\begin{itemize}
\item \cls{MatterPowerSpectrum}: $\sigma_8$, $\sigma_{12}$ and $A_s$; the linear $P(k)$ of the cold or total matter; the
      no-wiggle and BAO-smeared spectra; the nonlinear boost and $P_{\rm NL}$ of the cold or total matter, where the boost
      $Q$ multiplies the BAO-smeared spectrum
      $P_{\rm sm}=P_LG+P_{\rm nw}(1-G)$, $G=e^{-k^2\sigma_\ast^2/2}$, $\sigma_\ast^2=\int P_L\,\dd k/(3\pi^2)$, and the
      neutrinos enter as $P_{\rm tot}=(1-f_\nu)^2P_{\rm cold}+[P_{L,\rm tot}-(1-f_\nu)^2P_{L,\rm cold}]$; the baryonic
      suppression $S(k)=P_{\rm hydro}/P_{\rm grav}$ of a baryonification model (seven parameters: $M_c$, $\eta$, $\beta$,
      $M_{1,z_0,\rm cen}$, $\vartheta_{\rm out}$, $\vartheta_{\rm inn}$, $M_{\rm inn}$, or five for Burger et al.).
\item \cls{MatterBispectrum}: the baryonic suppression of the matter bispectrum for triangles $(k_1,k_2,k_3)$.
\item \cls{LbiasExpansion}: the 15 spectra $P_{ij}$ of the operators $(1,\delta,\delta^2,s^2,\nabla^2\delta)$ of the
      second order Lagrangian bias expansion (nonlinear, from the simulations, and LPT) for $k\le0.75\hMpc$, and the galaxy
      spectra $P_{gg}=\sum_{i\le j}(2-\delta_{ij})b_ib_jP_{ij}$, $P_{gm}=\sum_ib_iP_{0i}$ with $b_0=1$ and Lagrangian
      $(b_1,b_2,b_{s^2},b_{\nabla^2})$.
\item \cls{LbiasExpansionRSD}: the multipoles $\ell=0,2,4$ of the 15 spectra in redshift space ($k\le0.64\hMpc$), the galaxy
      $P(k,\mu)$ with fingers of god $\big[(1-f_{\rm sat})+f_{\rm sat}\lambda^2/(\lambda^2+k^2\mu^2)\big]^2$ and noise
      $(\epsilon_1+\epsilon_2k^2+\epsilon_3k^2\mu^2)/\bar n$, and its multipoles.
\item \nsp{COEVOLUTION}: the coevolution relations of the bias parameters of haloes and galaxies \cite{Zennaro2022}.
\end{itemize}

\begin{implementation}
The ReLU networks are evaluated with \cls{NeuralNetwork} (\cref{sec:fit-nn}), the networks of Burger et al.\ (CosmoPower
activation $(\beta+\mathrm{sigmoid}(\alpha x)(1-\beta))x$, with $\alpha$, $\beta$ per neuron) with a small evaluator of their
own; baccoemu evaluates in single precision, the port in double. Everything around the networks follows the Python package,
including its interpolations (linear in $\ln k$, or the not-a-knot cubic spline of \code{scipy.interpolate.interp1d} for the
bias spectra), the patches of the LPT terms in redshift space, and the projection of $P(k,\mu)$ on multipoles (a quartic fit
in $\mu$ on $\mu=0,0.01,\ldots,0.99$ integrated with an 8-point Gauss--Legendre rule). Parameters outside the training
ranges throw \cls{EmulationRangeError} naming the parameter.

The trained networks are not distributed with \fml{}: \file{Emulator/Bacco/get_data.sh} downloads them from bacco.dipc.org
(pinned SHA-256 sums; \code{--matter} for the matter emulators only, $\sim$10\,MB, otherwise $\sim$105\,MB) and converts
them with \file{export_bacco.py} (numpy and h5py) to \file{data/}: binary single-precision networks, 9\,MB for the matter
emulators and 89\,MB for the bias expansion. Without them the constructors throw \cls{MissingDataError}, which says how to
get them. Not ported: the BAO smearing via a DST of the linear spectrum (used by baccoemu only when both the smeared-BAO
and the no-wiggle emulators are off) and baccoemu's late-time distance helpers.
\end{implementation}

\begin{cppcode}[baccoemu: nonlinear and galaxy spectra]
#include <FML/Emulator/Bacco/Bacco.h>
using namespace FML::EMULATOR::BACCO;
Parameters p;
p.Omega_cold = 0.31; p.Omega_b = 0.049; p.sigma8_cold = 0.81;   // or p.Omega_m, p.A_s
p.h = 0.67; p.n_s = 0.965; p.M_nu = 0.06; p.w0 = -1.0; p.wa = 0.0;
MatterPowerSpectrum::Options opt;
opt.nonlinear_model = "Arico2023";
MatterPowerSpectrum emu(opt);
auto pnl = emu.nonlinear_pk(p, a, k, /*cold=*/false);     // total matter, (Mpc/h)^3
LbiasExpansion lb;
auto g = lb.galaxy_real_pk({b1, b2, bs2, bl}, p, a, k);    // g.auto_pk, g.cross_pk
\end{cppcode}

\begin{validation}[Validation: baccoemu port]
\file{Emulator/Bacco/tests} compares 251 cases with the Python package (every function, both nonlinear and both baryonic
models, cold and total matter with neutrinos, given linear spectra, the bias expansion in real and redshift space with and
without noise, and 18 out-of-range inputs, which name the same parameter). With the networks evaluated in double precision
in Python too the agreement is $\sim10^{-15}$; baccoemu's own single-precision evaluation moves its results by up to
$1.5\times10^{-6}$. This measures the port, not the emulators' accuracy against simulations, which is given in the papers.
\end{validation}

\section{Modified-gravity boosts}\label{sec:fit-mg}

\paragraph*{Hu--Sawicki $f(R)$.} \file{FofrFittingFunction.h} has the fits of Winther et al.\ \cite{Winther2019} to the
enhancement $B=P_{f(R)}/P_{\Lambda\rm CDM}$ for $n=1$ with the same initial conditions: \code{nonlinear_boost(z, fR0, k)}
(calibrated on the ELEPHANT simulations for $10^{-7}\le|f_{R0}|\le10^{-4}$, $0\le z\le3$, $0.01\le k\le10\hMpc$; three fits
around $|f_{R0}|=5\times10^{-6}$, $10^{-5}$, $5\times10^{-5}$ interpolated in $|f_{R0}|$; $\lesssim1\%$ on most scales, 2--3\%
at $k<10\hMpc$ out of sample) and \code{linear_boost} (one fit with the scaling $k\to k\sqrt{|f_{R0}|/10^{-5}}$). Outside the
ranges the value at the closest point is returned and the boost is never below 1. Equal to the original codes to
$10^{-13}$.

\paragraph*{Brans--Dicke.} For \code{model = jbd} the prescription of Sletmoen and Winther \cite{Sletmoen2024} maps the
Brans--Dicke universe to a GR universe with the same physical densities $\omega_b$, $\omega_c$, $n_s$ and $\sigma_8$ today
and $h_{\rm GR}=h_{\rm BD}\sqrt{\phi_{\rm ini}}$, and
\begin{equation}
P_{\rm BD}(k)=P^{\rm lin}_{\rm BD}(k)\,\frac{P_{\rm GR}(k')}{P^{\rm lin}_{\rm GR}(k')},\qquad \frac{k'}{h_{\rm GR}}=\frac{k}{h_{\rm BD}},
\end{equation}
(the same $k$ in $\hMpc$ of each universe), with the GR boost from \param{jbd_boost_gr} (\code{ee2} as in the paper, or
HMcode or Halofit). Accurate to 1\% against simulations for $\omega\ge100$, $z\le3$, $k\le1\hMpc$ \cite{Sletmoen2024}.
Against the paper's \file{bd.py} (hi\_class + EuclidEmulator2): $h_{\rm GR}$ to $10^{-7}$, $A_{s,\rm GR}$ to $2\times10^{-4}$
and the boost to $2\times10^{-4}$ for $z=0$--3 and $0.009\le k\le9\hMpc$.

\section{Eisenstein--Hu transfer functions}\label{sec:fit-eh}

\file{EisensteinHu.h} has the fitting formulae of Eisenstein and Hu \cite{EisensteinHu1998} for cold dark matter and baryons
(no neutrinos), with the baryon oscillations (their eqs.~16--24) and the zero-baryon-oscillation shape (eqs.~26--31), and the
shapes $k^{n_s}T^2(k)$. It is the one copy in the library, used for smooth shapes and extrapolations (the solver beyond
\param{pert_k_max}, HMcode-2020, the no-wiggle spectra, the field-level code), and equal to colossus' \code{eisenstein98}
and \code{eisenstein98_zb} to rounding.

\section{Neural-network emulators}\label{sec:fit-nn}

\file{NeuralNetwork.h} is a dependency-free fully connected network for emulators. It computes
\begin{equation}
\bm z=(\bm x-\bm\mu_x)/\bm s_x,\qquad \bm a_{\ell+1}=f(W_\ell\bm a_\ell+\bm b_\ell),\qquad
\bm y=(W_L\bm a_L+\bm b_L)\odot\bm s_y+\bm\mu_y,
\end{equation}
with the activations identity, ReLU, tanh, sigmoid, GELU (exact erf), SiLU, ELU, leaky ReLU and softplus, and the losses MSE,
L1 and Huber (as PyTorch). It does two things:
\begin{itemize}
\item \emph{Evaluate a network trained in PyTorch.} \file{export_lightning_emulator.py} exports a Lightning checkpoint of the
      Sesame pipeline's \file{Emulator.py} to a plain-text file; the C++ evaluation agrees with torch in double precision to
      $\sim10^{-16}$.
\item \emph{Train in C++}, with the same recipe and defaults as \file{Emulator.py}: Adam \cite{KingmaBa2014} (L2 weight decay,
      $\beta_1=0.9$, $\beta_2=0.999$), Glorot initialisation \cite{Glorot2010}, gradient clipping (norm 0.5), ReduceLROnPlateau
      (patience 20, factor 0.1), early stopping (patience 30) and stochastic weight averaging \cite{Izmailov2018}.
\end{itemize}
\code{jacobian} gives $\partial y_k/\partial x_j$ by back-propagation, for Fisher matrices or Hamiltonian Monte Carlo
(\cref{ch:sampling-fitting}); \code{evaluate} runs many points with OpenMP; \code{save} writes 17 significant digits so a
file read back gives a bit-identical network. The training points are best chosen with a Latin hypercube
\cite{McKay1979} (\file{LatinHyperCube}). The example emulates the linear power spectrum in $(\Om,\sigma_8,n_s)$: it trains in
a few seconds and reaches a mean $|\Delta P/P|=7\times10^{-4}$ on 200 test cosmologies.

\begin{cppcode}[Train, use and save an emulator]
#include <FML/Emulator/NeuralNetwork/NeuralNetwork.h>
using namespace FML::EMULATOR::NEURALNETWORK;
TrainingOptions opt;
opt.hidden_layers = {64, 64};
opt.activation = Activation::gelu;
TrainingHistory history;
NeuralNetwork net = train(X_train, Y_train, X_val, Y_val, opt, &history);
std::vector<double> y = net(x);              // one point
Matrix J = net.jacobian(x);                  // dy_k / dx_j
net.save("emulator.txt");
auto net2 = NeuralNetwork::load("emulator.txt");   // also networks exported from PyTorch
\end{cppcode}
